\section{Fr\'echet dispersion as residual link energy}
\label{sec:frechet}

\begin{definition}[Staple dispersion]
For $m$ complementary transports in one transport bitorsor, define
\begin{align}
 \Phi(C)&=\sum_{i=1}^m d(C,S_i)^2,\\
 D(S_1,\ldots,S_m)&=\inf_C\Phi(C).
\end{align}
A minimizer, when chosen, is denoted $C_*$.  The quantity $D$ is the squared
Fr\'echet dispersion of the star.
\end{definition}

The definition is independent of a trivialization.  A frame change sends
every $C$ and $S_i$ to $g_y(\cdot)g_x^{-1}$, an isometry of the bi-invariant
metric.

\begin{theorem}[Residual-energy identity]
\label{thm:residual}
For the geodesic-quadratic plaquette energy incident on $\ell$,
\[
  \mathcal E_\ell(U_\ell)=\sum_{i=1}^m d(I,U_\ell S_i)^2,
\]
one has the exact identity
\[
  \mathcal E_\ell(U_\ell)=\Phi(U_\ell^{-1}),
  \qquad
  \inf_{U_\ell\in\SU}\mathcal E_\ell(U_\ell)=D(S_1,\ldots,S_m).
\]
In particular, $\mathcal E_\ell(U_\ell)\ge D$ for every central link.
\end{theorem}

The proof is the one-line bi-invariance calculation
$d(I,US)=d(U^{-1},S)$, but the interpretation is useful.  $D$ is not simply
``large plaquette curvature.''  It is the part of the incident energy that
cannot be removed by any value of the shared link.  A star with $D=0$ is
perfectly coherent: all complementary transports agree.  Positive $D$
measures relational frustration among them.

\begin{corollary}[Diameter witness]
If no geodesic ball of radius $\eta$ contains all $S_i$, then $D>\eta^2$.
Equivalently, $D\le\eta^2$ implies that the ball of radius $\eta$ centered at
any minimizer contains every staple.
\end{corollary}

This witness is intentionally one-sided.  A large containing radius does not
force $D$ to be extensive in $m$ without additional separation information.
The pairwise formulation below captures that information more symmetrically.
